\begin{tcolorbox}[
    colback=SkyBlue!5,
    colframe=blue!20,
    boxrule=1pt,
    arc=2pt,
    width=\textwidth,
    boxsep=5pt,
    left=5pt,
    right=5pt,
    top=5pt,
    bottom=5pt
]

\noindent
\begin{minipage}[b]{0.62\textwidth}
\begin{example} \label{example:R+RR}
Unfortunately, \textbf{Algorithm P} can have very high variance even for a trivial program.
Consider this simple program with three events:
 \begin{equation*}
  \tag{\textsc{r+rr}}
  \label{ex:r+rr}
  \inarr{ \textcolor{Red}{\readInst{a_1}{y}} }{
\textcolor{blue}{\readInst{a_2}{x}} \\
\textcolor{blue}{\readInst{a_3}{x}}
}   
\end{equation*}
\cref{meta-graph-r+rr} shows the DAG $\mathcal{T}(\text{\textsc{r+rr}})$ of the program. Although $\mathcal{T}(\text{\textsc{r+rr}})$ has exactly one sink, there are three paths to reach the sink, and running \textbf{Algorithm P} on each path returns a different estimation.
\end{example}
\end{minipage}%
\hfill
\begin{minipage}[b]{0.40\textwidth}
\centering
\vspace{0pt}
\centering
\begin{tikzpicture}[
    scale=0.7,
    transform shape,
    >=stealth,
    node distance=2.2cm and 2.5cm,
    hypernode/.style={draw, rectangle, rounded corners=4pt, very thick, inner sep=3pt, font=\small\bfseries},
    event/.style={draw=none, font=\normalsize},
    rf/.style={->, thick, dashed, green!60!black},
    co/.style={->, thick, dotted, red!80!black},
    graph box/.style={draw, rectangle, rounded corners=4pt, thick, minimum width=1.4cm, minimum height=1cm, inner sep=0.1cm, minimum size=0cm},
    edge label/.style={midway, very thick, fill=white, font=\footnotesize, inner sep=1pt},
    edge-t1/.style={->, very thick, Red},
    edge-t2/.style={->, very thick, blue},
]

% Top layer: G0
\node[hypernode] (G0) at (0,0) {$G_0$};

\node[hypernode] (G1) at (-1,-1) {$G_1$};

\node[hypernode] (G2) at (1,-1) {$G_2$};

\node[hypernode] (G3) at (0,-2) {$G_3$};

\node[hypernode] (G4) at (2,-2) {$G_4$};

% G5 - Concrete graph
\begin{scope}[xshift=2cm, yshift=-3.2cm, scale=0.5, local bounding box=G5box]
  \node[event] (g5i) {$\init$};
  \node[event, below=0.7cm of g5i, xshift=-1.0cm] (g5r1) {$\opr(y)$};
  \node[event, below=0.7cm of g5i, xshift=1.0cm] (g5r2) {$\opr(x)$};
  \node[event, below=0.4cm of g5r2] (g5r3) {$\opr(x)$};
  \draw[->] (g5i) -- (g5r1);
  \draw[->] (g5i) -- (g5r2);
  \draw[->] (g5r2) -- (g5r3);
\end{scope}
\node[graph box, fit=(G5box)] (G5) {};

% Abstract hyperedges
\draw[edge-t1] (G0) -- (G1);
\draw[edge-t2] (G0) -- (G2);

\draw[edge-t2] (G1) -- (G3);
\draw[edge-t1] (G2) -- (G3);

% Edges from G2
\draw[edge-t2] (G2) -- (G4);

% Edges to G5 - Concrete graph
\draw[edge-t2] (G3) -- (G5);
\draw[edge-t1] (G4) -- (G5);
\end{tikzpicture}
\captionof{figure}{The transition system $\mathcal{T}(\ref{ex:r+rr})$}
\label{meta-graph-r+rr}


% \begin{figure}[htbp]
% \centering
% \resizebox{0.6\textwidth}{!}{ 
% \begin{tikzpicture}[
%     scale=0.7,
%     transform shape,
%     >=stealth,
%     node distance=2.2cm and 2.5cm,
%     hypernode/.style={draw, rectangle, rounded corners=4pt, very thick, inner sep=3pt, font=\small\bfseries},
%     edge label/.style={midway, very thick, fill=white, font=\footnotesize, inner sep=1pt},
%     edge-t1/.style={->, very thick, Red},
%     edge-t2/.style={->, very thick, blue},
%     edge-t3/.style={->, very thick, Dandelion!60!black}
% ]

% node[hypernode] (G0) at (0,0) {G0};


% \begin{scope}[xshift=0cm, yshift=0cm, local bounding box=G0box]
%   \node[event] (g0i) {$\init$};
% \end{scope}
% \node[graph box, fit=(G0box)] (G0) {};

% % G1
% \begin{scope}[xshift=-4cm, yshift=-2cm, local bounding box=G1box]
%   \node[event] (g1i) {$\init$};
%   \node[event, below=0.7cm of g1i, xshift=-1.0cm] (g1r1) {$\opr(y)$};
%   \draw[->] (g1i) -- (g1r1);
% \end{scope}
% \node[graph box, fit=(G1box)] (G1) {};

% % G2
% \begin{scope}[xshift=4cm, yshift=-2cm, local bounding box=G2box]
%   \node[event] (g2i) {$\init$};
%   \node[event, below=0.7cm of g2i, xshift=1.0cm]  (g2r2) {$\opr(x)$};
%   \draw[->] (g2i) -- (g2r2);
% \end{scope}
% \node[graph box, fit=(G2box)] (G2) {};


% % G4
% \begin{scope}[xshift=0cm, yshift=-5cm, local bounding box=G4box]
%   \node[event] (g4i) {$\init$};
%   \node[event, below=0.7cm of g4i, xshift=-1.0cm] (g4r1) {$\opr(y)$};
%   \node[event, below=0.7cm of g4i, xshift=1.0cm]  (g4r2) {$\opr(x)$};
%   \draw[->] (g4i) -- (g4r1);
%   \draw[->] (g4i) -- (g4r2);
% \end{scope}
% \node[graph box, fit=(G4box)] (G4) {};

% % G5
% \begin{scope}[xshift=7cm, yshift=-5cm, local bounding box=G5box]
%   \node[event] (g5i) {$\init$};
%   \node[event, below=0.7cm of g5i, xshift=1.0cm] (g5r1) {$\opr(x)$};
%   \node[event, below=0.4cm of g5r1] (g5l2) {$\opr(x)$};
%   \draw[->] (g5i) -- (g5r1);
%   \draw[->] (g5r1) -- (g5l2);

% \end{scope}
% \node[graph box, fit=(G5box)] (G5) {};
% % G8
% \begin{scope}[xshift=4cm, yshift=-8cm, local bounding box=G8box]
%   \node[event] (g8i) {$\init$};
%   \node[event, below=0.7cm of g8i, xshift=-1.0cm] (g8r1) {$\opr(y)$};
%   \node[event, below=0.7cm of g8i, xshift=1.0cm] (g8r2) {$\opr(x)$};
%   \node[event, below=0.4cm of g8r2] (g8l3) {$\opr(x)$};
%   \draw[->] (g8i) -- (g8r1);
%   \draw[->] (g8i) -- (g8r2);
%   \draw[->] (g8r2) -- (g8l3);
% \end{scope}
% \node[graph box, fit=(G8box)] (G8) {};




% \draw[hyper edge] (G0) -- node[midway, above left, font=\small] {$t_1$} (G1);
% \draw[hyper edge] (G0) -- node[midway, above, font=\small] {$t_2$} (G2);
% \draw[hyper edge] (G1) -- node[midway, above, font=\small] {$t_2$} (G4);
% \draw[hyper edge] (G2) -- node[midway, above, font=\small] {$t_1$} (G4);
% \draw[hyper edge] (G2) -- node[midway, above, font=\small] {$t_2$} (G5);

% \draw[hyper edge] (G4) -- node[midway, above, font=\small] {$t_2$} (G8);
% \draw[hyper edge] (G5) -- node[midway, above, font=\small] {$t_1$} (G8);
% \end{tikzpicture}
% }
% \caption{Transition system $\mathcal{T}(\text{\textsc{r+rr}})$}
% \label{meta-graph-r+rr}
% \end{figure}
\end{minipage}%

For example, the leftmost path—where $t_1$ is scheduled first—returns $0.5$. The middle path—where $t_1$ is scheduled second—returns $1$. The rightmost path—where $t_1$ is scheduled last—returns $2$. Intuitively, the path in which all events of $t_2$ are scheduled before $t_1$ has the lowest probability of being chosen and therefore yields the largest estimate. 

We extend this to the case where $t_2$ performs $n$ reads from $x$ (with $n>2$). There is again exactly one sink, but $n{+}1$ root-to-sink paths, indexed by $k\in[0,n]$ as the number of $t_2$ events before $t_1$. The path with index $k$ is taken with probability $2^{-k-1}$ if $k\neq n$ and $2^{-n}$ if $k=n$. The estimate returned by \textbf{Algorithm P} is 
\[
r_0 = 2^{1-n}, \quad r_k = 2^{2k-n} \mbox{ for }0< k <n, \mbox{ and }r_n = 2^{n-1}
\]
%
The variance of \textbf{Algorithm P} on this program is $\operatorname{Var}(r)=\frac{9}{28}\,2^{n}+\frac{10}{7}\,2^{-2n}-1$, which is exponential in $n$. Since the variance is large in this very simple setting, scaling to more complex programs is not promising using this approach.
\qed
\end{tcolorbox}